% !TEX root = ../../main.tex
% C3.2 - Person detection (translated). Matches config/ultralytics_yolo_detector.json (conf 0.1, iou 0.7, max 300).
\section{Person Detection}
\label{sec:3.2}

\subsection{Object Detection}
\label{sec:3.2.1}

Object detection locates the objects in an image and determines their class. For each object, the model returns a rectangular \emph{bounding box}, usually given by the coordinates of two corners $(x_1, y_1, x_2, y_2)$, the \emph{class} of the object and a \emph{confidence score}. Two boxes are compared with the Intersection over Union (IoU):
\begin{equation}
    \mathrm{IoU}(A, B) = \frac{|A \cap B|}{|A \cup B|},
    \label{eq:iou}
\end{equation}
which ranges from 0 when the boxes do not overlap to 1 when they coincide. IoU is used to remove duplicate predictions and to link detections across frames in tracking (Section~\ref{sec:3.3}).

\subsection{Two-Stage and One-Stage Detectors}
\label{sec:3.2.2}

A \emph{two-stage} model such as Faster R-CNN~\cite{ren2017fasterrcnn} first uses a region proposal network to generate regions likely to contain objects, then classifies each region and refines its box; this is accurate but costly because many proposals must be processed. A \emph{one-stage} model such as YOLO (You Only Look Once)~\cite{redmon2016yolo} treats detection as a single regression problem, predicting the boxes and classes of all objects in one pass over the image, which is considerably faster and suits processing many video frames.

\subsection{The YOLO Family}
\label{sec:3.2.3}

Modern YOLO models have three parts: a feature extractor (backbone) at several levels of detail, a feature aggregator (neck) at several scales so that both large objects close to the camera and small distant ones can be detected, and a prediction head for boxes, classes and confidence. YOLO11 was released by Ultralytics in 2024 in five sizes (n, s, m, l, x) that trade speed for accuracy~\cite{ultralytics2024yolo11docs}. The models are pre-trained on COCO~\cite{lin2014coco}, which has 80 common object classes including \emph{person}, so they can detect people directly without further training.

\subsection{Post-Processing: Thresholding and Duplicate Removal}
\label{sec:3.2.4}

The raw output has many low-confidence and overlapping boxes. Boxes below a confidence threshold are discarded, and Non-Maximum Suppression (NMS) keeps the most confident box and removes same-class boxes overlapping it beyond an IoU threshold. A lower confidence threshold misses fewer people but adds false detections; a lower IoU threshold removes overlaps more aggressively but may merge people standing close together.

\subsection{Role of Detection in the System}
\label{sec:3.2.5}

The detector is applied to the sampled frames (Section~\ref{sec:3.4}) and keeps only the \emph{person} class, with a confidence threshold of 0.1, an NMS IoU threshold of 0.7 and at most 300 detections per frame. The low threshold is deliberate: ByteTrack needs low-confidence detections, often partially occluded people, to keep tracks alive (Section~\ref{sec:3.3}); they never start tracks, and a track needs two observations to be confirmed, so most transient false detections are removed by tracking. The reasons for choosing YOLO11n are given in Section~\ref{sec:6.1.4}.
